\documentclass[a4paper,11pt]{article}

\usepackage[utf8]{inputenc}
\usepackage[T1]{fontenc}
\usepackage[ngerman]{babel}
\usepackage{amsmath}
\usepackage{amssymb}
\usepackage{amsthm}
\usepackage{xcolor}
\usepackage[most]{tcolorbox}
\usepackage{titlesec}
\usepackage{enumitem}
\usepackage{tikz}
\usepackage{geometry}
\geometry{a4paper,margin=2.1cm}

% ---------- Dark Mode ----------
\definecolor{bgdark}{HTML}{1E1E1E}
\definecolor{panel}{HTML}{2A2A2A}
\definecolor{fgtext}{HTML}{E6E6E6}
\definecolor{accent}{HTML}{89B4FA}  % blau  -> Definition
\definecolor{green}{HTML}{A6E3A1}    % grün  -> Intuition
\definecolor{yellow}{HTML}{F9E2AF}   % gelb  -> Satz/Formel
\definecolor{red}{HTML}{F38BA8}      % rot   -> Stolperfalle
\definecolor{muted}{HTML}{A6ADC8}
\definecolor{rulecol}{HTML}{6A6A6A}
\pagecolor{bgdark}
\color{fgtext}
\renewcommand{\normalcolor}{\color{fgtext}}   % Überschriften/TOC hell

\usepackage[colorlinks=true,linkcolor=accent,urlcolor=accent]{hyperref}

\setlength{\parindent}{0pt}
\setlength{\parskip}{4pt}

% ---------- Section-Stil ----------
\titleformat{\section}{\Large\bfseries\color{accent}}{\thesection}{0.6em}{}
\titleformat{\subsection}{\large\bfseries\color{yellow}}{\thesubsection}{0.6em}{}
\titlespacing*{\section}{0pt}{10pt}{4pt}
\titlespacing*{\subsection}{0pt}{7pt}{2pt}

% ---------- Boxen ----------
\newtcolorbox{defi}[1]{colback=panel,colframe=accent,coltext=fgtext,coltitle=bgdark,
  colbacktitle=accent,fonttitle=\bfseries,title={#1},arc=2pt,boxrule=0.6pt,
  left=4pt,right=4pt,top=2pt,bottom=2pt,before skip=4pt,after skip=4pt}
\newtcolorbox{intu}[1]{colback=panel,colframe=green,coltext=fgtext,coltitle=bgdark,
  colbacktitle=green,fonttitle=\bfseries,title={#1},arc=2pt,boxrule=0.6pt,
  left=4pt,right=4pt,top=2pt,bottom=2pt,before skip=4pt,after skip=4pt}
\newtcolorbox{satz}[1]{colback=panel,colframe=yellow,coltext=fgtext,coltitle=bgdark,
  colbacktitle=yellow,fonttitle=\bfseries,title={#1},arc=2pt,boxrule=0.6pt,
  left=4pt,right=4pt,top=2pt,bottom=2pt,before skip=4pt,after skip=4pt}
\newtcolorbox{falle}[1]{colback=panel,colframe=red,coltext=fgtext,coltitle=bgdark,
  colbacktitle=red,fonttitle=\bfseries,title={#1},arc=2pt,boxrule=0.6pt,
  left=4pt,right=4pt,top=2pt,bottom=2pt,before skip=4pt,after skip=4pt}

\newcommand{\K}{\textcolor{green}{\ding}}
\setlist[itemize]{leftmargin=1.4em,itemsep=1pt,topsep=2pt}

\begin{document}

% ---------- Titel ----------
\begin{center}
{\Huge\bfseries\color{accent} Numerische Methoden}\\[4pt]
{\large\color{fgtext} Theorie-Lernskript -- alle Themen kompakt erklärt}\\[3pt]
{\small\color{muted} Grundlagen $\cdot$ Fehleranalyse $\cdot$ Nullstellen $\cdot$ Optimierung $\cdot$ DGL $\cdot$ Integration/Interpolation $\cdot$ Fourier $\cdot$ Modelltraining}
\end{center}
\vspace{2pt}
\noindent\textcolor{rulecol}{\rule{\linewidth}{0.4pt}}

{\small\color{muted} Legende der Boxen:\;
\textcolor{accent}{$\blacksquare$}~Definition \quad
\textcolor{green}{$\blacksquare$}~Intuition \quad
\textcolor{yellow}{$\blacksquare$}~Satz/Formel \quad
\textcolor{red}{$\blacksquare$}~Stolperfalle (richtige Form direkt gezeigt)}

\vspace{4pt}
{\color{muted}\hypersetup{linkcolor=accent}\tableofcontents}
\vspace{6pt}
\noindent\textcolor{rulecol}{\rule{\linewidth}{0.4pt}}

% =====================================================================
\section{Grundlagen \& Zahlendarstellung}

\begin{intu}{Warum Numerik?}
Viele Probleme haben \emph{keine} geschlossene Lösung (z.\,B.\ $\int e^{-x^2}dx$, $x=\cos x$,
allgemeine Nullstellen ab Grad~5) oder sind analytisch zu aufwendig. \textbf{Numerik} liefert
\emph{näherungsweise} Zahlenergebnisse mit endlich vielen Rechenschritten -- Preis: jedes Ergebnis
hat einen \emph{Fehler}, den man verstehen und kontrollieren muss.
\end{intu}

\subsection{Diskretisierung \& Brute Force}
Ein kontinuierliches Problem wird durch endlich viele Stützstellen ersetzt. Intervall $[a,b]$ in
$N$ Teile: \textbf{Schrittweite} $h=\frac{b-a}{N}$, Gitterpunkte $x_i=a+i\,h$ ($i=0,\dots,N$, also
$N{+}1$ Punkte). \emph{Grid Search / Brute Force}: alle Gitterpunkte durchprobieren und den besten
nach einem \emph{Fehlermaß} wählen.

\begin{satz}{Grid Search für $A\mathbf x=\mathbf c$}
Fehlermaß = \textbf{Residuum} $E(\mathbf x)=\lVert A\mathbf x-\mathbf c\rVert$ (z.\,B.\ Summe der
Beträge oder Quadrate). Über das Gitter minimieren. \emph{Trifft die exakte Lösung nur, wenn diese
zufällig auf einem Gitterpunkt liegt} -- sonst bleibt ein Diskretisierungsfehler. Verbesserung:
feineres Gitter oder lokale Verfeinerung um den besten Kandidaten.
\end{satz}

\subsection{Gleitkomma (floating point)}
\begin{defi}{Gleitkommazahl}
$x=(-1)^s\,m\cdot 2^{e}$ mit \textbf{Vorzeichen} $s$, \textbf{Mantisse} $m$ (signifikante Stellen),
\textbf{Exponent} $e$ (Skalierung, gespeichert mit \emph{Bias}: $e=E-\text{Bias}$).
\emph{Normalisiert}: $m=1{.}b_1b_2\dots b_t$ mit \emph{implizitem} führenden $1$-Bit (1 Gratis-Bit).
\end{defi}

\begin{satz}{Maschinenepsilon}
$\varepsilon_{\mathrm{mach}}=2^{-t}$ ($t$ = Anzahl gespeicherter Mantissen-Nachkommabits) ist der
Abstand von $1$ zur nächstgrößeren darstellbaren Zahl -- die \emph{kleinste relative Auflösung}.
Relativer Rundungsfehler stets $\le \tfrac12\varepsilon_{\mathrm{mach}}$.
\end{satz}

\begin{intu}{Trade-off Mantisse $\leftrightarrow$ Exponent (bei fester Bitzahl)}
\textbf{Mehr Mantissenbits} $\Rightarrow$ kleineres $\varepsilon$ $\Rightarrow$ feinere
\emph{Auflösung} (resolution). \textbf{Mehr Exponentenbits} $\Rightarrow$ größerer
\emph{Wertebereich} (dynamic range, größte/kleinste Beträge). Man tauscht also stets
\emph{Präzision} gegen \emph{Reichweite}.
\end{intu}

\subsection{Festkomma (fixed point)}
$f$ Nachkommabits, \textbf{Auflösung} $=2^{-f}$ (konstant). Speicherwert $=$ Zahl$\,\cdot 2^{f}$
(ganzzahlig), größte vorzeichenlose Zahl mit $k$ Gesamtbits $=\frac{2^{k}-1}{2^{f}}$.

\begin{falle}{Festkomma vs.\ Gleitkomma -- konstanter \emph{absoluter} vs.\ \emph{relativer} Abstand}
Festkomma: Nachbarzahlen haben \emph{überall denselben} Abstand $2^{-f}$ (gleichmäßiges Raster).
Gleitkomma: \emph{relativer} Abstand konstant ($\approx\varepsilon$), d.\,h.\ der \emph{absolute}
Abstand wächst mit der Größenordnung von $x$ (eng bei kleinen, weit bei großen Zahlen).
\end{falle}

\subsection{Rundung, Auslöschung \& Fehlerarten}
\begin{satz}{Auslöschung (catastrophic cancellation)}
Subtrahiert man zwei \emph{fast gleiche} Zahlen, löschen sich die führenden Stellen weg; übrig
bleiben wenige, schon verrauschte Stellen $\Rightarrow$ \emph{relativer} Fehler explodiert.
Beispiel: $(1+\delta)-1$ verliert $\delta$ komplett, falls $\delta<\tfrac12\varepsilon$
(dann Ergebnis $0$, Division dadurch $\to\infty$).
\end{satz}

\begin{intu}{Gegenmittel: umformen statt subtrahieren}
$\sqrt{x+1}-\sqrt{x}=\dfrac{1}{\sqrt{x+1}+\sqrt{x}}$ -- rechts wird \emph{addiert}, keine
Auslöschung mehr. Allgemein: kritische Differenzen algebraisch in Summen/Produkte umschreiben.
\end{intu}

\begin{defi}{Die drei Fehlerarten}
\begin{itemize}
\item \textbf{Modellfehler}: das Modell vereinfacht die Realität (z.\,B.\ Reibung/Luftwiderstand vernachlässigt).
\item \textbf{Abschneide-/Diskretisierungsfehler}: ein Grenzprozess wird abgebrochen
  (Reihe nach endlich vielen Gliedern; Ableitung $\to$ Differenzenquotient). Skaliert mit $h$.
\item \textbf{Rundungsfehler}: endliche Stellenzahl der Maschine (z.\,B.\ $\tfrac{10}{3}\to 3{,}333$).
\end{itemize}
\textbf{Absoluter} Fehler $|x-\tilde x|$; \textbf{relativer} $\frac{|x-\tilde x|}{|x|}$
(skaleninvariant, „gültige Stellen``; bei $x\approx0$ nur absolut sinnvoll).
\end{defi}

% =====================================================================
\section{Fehleranalyse: die vier Eigenschaften}

\begin{intu}{Hut \& Tilde}
$f$ = wahres Problem, $\hat f$ = der \emph{Algorithmus} (gestörte Rechenvorschrift),
$\tilde x$ = \emph{gestörte Eingabe}. Die vier Eigenschaften vergleichen diese vier Größen
paarweise.
\end{intu}

\begin{defi}{Kondition, Stabilität, Konsistenz, Konvergenz}
\begin{itemize}
\item \textbf{Kondition} $\kappa=|f(x)-f(\tilde x)|$: wie stark reagiert das \emph{Problem} auf
  Eingabestörung. Klein $=$ gut, groß $=$ schlecht konditioniert. \emph{Eigenschaft des Problems.}
\item \textbf{Stabilität} $|\hat f(\tilde x)-\hat f(x)|\le s\,|\tilde x-x|$: wie stark verstärkt der
  \emph{Algorithmus} Störungen. \emph{Eigenschaft des Verfahrens.}
\item \textbf{Konsistenz} $|\hat f(x)-f(x)|\le c$: wie gut nähert der Algorithmus bei \emph{exakter}
  Eingabe.
\item \textbf{Konvergenz} $|\hat f(\tilde x)-f(x)|\to 0$: Gesamtfehler aus gestörter Eingabe.
\end{itemize}
\end{defi}

\begin{satz}{Kernzusammenhang}
\textbf{Konsistenz $+$ Stabilität $\Rightarrow$ Konvergenz.} Quantifizierung über $|f'(x)|$:
$\kappa\approx|f'(x)|\,\delta x$ -- dort, wo $f$ \emph{flach} ist (kleine Steigung), gut
konditioniert; wo $f$ \emph{steil} ist, schlecht.
\end{satz}

\begin{falle}{Zwei Verwechslungen, die Punkte kosten}
1) \textbf{Kondition $\ne$ Stabilität}: Kondition gehört zum \emph{Problem}, Stabilität zum
\emph{Algorithmus}. Ein stabiler Algorithmus kann ein \emph{schlecht konditioniertes} Problem
\textbf{nicht} retten -- er erzeugt nur keine \emph{zusätzlichen} Fehler.\\
2) \textbf{Kleineres $h$ ist \emph{nicht} automatisch besser}: feineres $h$ senkt den
Konsistenz-/Diskretisierungsfehler, aber unterhalb eines optimalen $h$ \emph{dominiert der
Rundungsfehler} (und Eingangsrauschen) -- der Gesamtfehler steigt wieder. Es gibt ein optimales $h$.
\end{falle}

% =====================================================================
\section{Nullstellen: Newton- \& Sekantenverfahren}

\begin{satz}{Newton-Verfahren -- Herleitung aus Taylor}
Taylor 1.\ Ordnung: $f(x)\approx f(x_n)+f'(x_n)(x-x_n)$. Diese \emph{Tangente} $=0$ setzen und nach
$x$ lösen:
\[
  \boxed{\,x_{n+1}=x_n-\frac{f(x_n)}{f'(x_n)}\,}
\]
Newton iteriert also die Nullstelle der Tangente.
\end{satz}

\begin{defi}{Konvergenz von Newton}
\textbf{Quadratische Konvergenz} $|e_{n+1}|\le C|e_n|^2$ nahe einer \emph{einfachen} Wurzel
$x^\star$: die Zahl korrekter Stellen \emph{verdoppelt} sich pro Schritt
($10^{-1}\!\to\!10^{-2}\!\to\!10^{-4}\!\to\!10^{-8}$).
\textbf{Konvergenzkriterium}: $x_0$ nah genug an einfacher Wurzel, $f\in C^2$, $f'(x^\star)\neq 0$.
\end{defi}

\begin{falle}{Newton-Versagensfälle}
\begin{itemize}
\item $f'(x_n)=0$: waagrechte Tangente $\Rightarrow$ \emph{Division durch $0$} (Verfahren bricht ab).
\item Schlechter Startwert $\Rightarrow$ Divergenz oder Sprung ins falsche Becken.
\item \emph{Zyklus}: Iteration pendelt periodisch und konvergiert nie.
\end{itemize}
\end{falle}

\begin{satz}{Finite Differenz \& Sekantenverfahren}
Ist $f'$ unbekannt/teuer: \textbf{Vorwärtsdifferenz}
$f'(x_0)\approx\frac{f(x_0+h)-f(x_0)}{h}$, Fehler $\mathcal O(h)$ (Taylor-Restterm $\tfrac h2 f''$).
Einsetzen $\to$ \textbf{Sekantenverfahren}
\[
  x_{n+1}=x_n-f(x_n)\,\frac{x_n-x_{n-1}}{f(x_n)-f(x_{n-1})}.
\]
\emph{Ableitungsfrei} (Ordnung $\approx1{,}618$). Vorteil: kein $f'$ nötig. Nachteil: langsamer als
Newton, braucht zwei Startwerte, Nenner kann $\approx0$ werden.
\end{satz}

% =====================================================================
\section{Optimierung}

\begin{defi}{Lokal vs.\ global, Konvexität}
\textbf{Lokales} Minimum: kleinster Wert in einer \emph{Umgebung}; \textbf{globales}: kleinster Wert
auf der \emph{ganzen} Domäne. Lokale Verfahren folgen nur lokaler Information und laufen ins
\emph{nächste} Tal -- bei \emph{multimodalen} Landschaften (viele Minima, z.\,B.\ Shekel-Foxholes)
verfehlen sie das globale.
\end{defi}

\begin{intu}{Konvexität als Glücksfall}
Ist $f$ \emph{konvex} ($f''\ge 0$ überall), gibt es nur \emph{ein} Minimum -- dann ist jedes lokale
zugleich global, lokale Verfahren genügen. Schwierig wird es erst im \emph{nicht-konvexen} Fall.
\end{intu}

\begin{satz}{Newton fürs Optimum \& Gradientenabstieg}
Optimum: $f'(x)=0$. Newton auf $f'$:
\[
  x_{n+1}=x_n-\frac{f'(x_n)}{|f''(x_n)|}.
\]
Der \emph{Betrag} $|f''|$ erzwingt einen Schritt \emph{immer bergab} (auch in konkaven Regionen
$f''<0$, wo reines Newton bergauf zu einem Maximum liefe). Klassifikation via 2.\ Ableitung:
$f''>0\Rightarrow$ Minimum, $f''<0\Rightarrow$ Maximum.\\
\textbf{Gradientenabstieg}: $x_{n+1}=x_n-\eta f'(x_n)$. Lernrate $\eta$ zu \emph{groß}
$\Rightarrow$ Überschießen/Oszillation/Divergenz; zu \emph{klein} $\Rightarrow$ sehr langsam.
\end{satz}

\begin{satz}{Globale Strategien}
\begin{itemize}
\item \textbf{Random Restarts}: viele zufällige Starts; Erfolgswahrscheinlichkeit nach $K$ Starts
  $P=1-(1-p)^K$ ($p$ = Trefferchance pro Start).
\item \textbf{Simulated Annealing}: akzeptiert Verschlechterung $\Delta E>0$ mit
  $e^{-\Delta E/T}$. Hohe Temperatur $T$ = viel Exploration (verlässt lokale Minima), Abkühlen
  $T\downarrow$ = Exploitation.
\item \textbf{Basin Hopping}: zufälliger Sprung $+$ lokale Minimierung, springt zwischen „Becken``.
\end{itemize}
\end{satz}

\begin{falle}{Raue Landschaft \& Glättung -- Wahl von $h$}
Ein hochfrequenter Störterm wie $\tfrac12\sin(4\pi x)$ erzeugt viele Mikro-Minima; lokale
Optimierer bleiben im ersten stecken. \textbf{Glätten} (Fenster-Mittel/Integration über eine
Periode) bügelt die Oszillation aus. Wähle $h\approx$ \emph{Periode} des Störterms:
$\sin(4\pi x)$ hat Periode $\frac{2\pi}{4\pi}=\tfrac12$ $\Rightarrow$ $h\approx\tfrac12$.
\end{falle}

% =====================================================================
\section{Differenzialgleichungen (Anfangswertprobleme)}

\begin{defi}{Euler-Verfahren}
AWP $y'=f(t,y),\ y(t_0)=y_0$.
\textbf{Explizit}: $y_{n+1}=y_n+h\,f(t_n,y_n)$ (direkt, billig).
\textbf{Implizit}: $y_{n+1}=y_n+h\,f(t_{n+1},y_{n+1})$ (Gleichung lösen, teurer, viel stabiler).
\end{defi}

\begin{satz}{Runge-Kutta \& Butcher-Tableau}
RK kombiniert mehrere Steigungs-Auswertungen. \textbf{RK4}:
\[
  y_{n+1}=y_n+\tfrac h6(k_1+2k_2+2k_3+k_4),
\]
\[
  k_1=f(t_n,y_n),\;\; k_2=f(t_n{+}\tfrac h2,\,y_n{+}\tfrac h2 k_1),\;\;
  k_3=f(t_n{+}\tfrac h2,\,y_n{+}\tfrac h2 k_2),\;\; k_4=f(t_n{+}h,\,y_n{+}h\,k_3).
\]
\textbf{Butcher-Tableau} $\begin{array}{c|c}c&A\\\hline&b^{\!\top}\end{array}$:
Knoten $c_i$ (Zeitpunkte), Gewichte $b_i$ (Endmischung), Matrix $a_{ij}$ (Stufenkopplung).
\emph{Explizit} $\Leftrightarrow$ $A$ streng untere Dreiecksmatrix. Ordnungsbedingungen u.\,a.\
$\sum_i b_i=1$, $\sum_i b_i c_i=\tfrac12$.
\end{satz}

\begin{intu}{Ordnung, Stabilität, Steifheit}
\textbf{Konvergenzordnung} $p$: globaler Fehler $\mathcal O(h^p)$ (Euler $p{=}1$, RK4 $p{=}4$). RK4
ist bei gleichem $h$ \emph{drastisch} genauer (4 Auswertungen, dafür viel größere Schritte möglich).
\textbf{Stabilität} (explizit Euler, $y'=-\lambda y$): stabil nur für $0<h<\tfrac{2}{\lambda}$.
\textbf{Steifheit}: stark unterschiedliche Zeitskalen zwingen explizite Verfahren zu winzigen $h$
$\Rightarrow$ implizite Verfahren nötig.
\end{intu}

% =====================================================================
\section{Numerische Integration \& Interpolation}

\begin{defi}{Newton-Cotes-Idee}
Funktion durch ein Interpolationspolynom über äquidistante Stützstellen ersetzen und \emph{dieses}
exakt integrieren. \textbf{Exaktheitsgrad} = höchster Polynomgrad, den die Regel noch exakt
integriert.
\end{defi}

\begin{satz}{Die Quadraturregeln (auf $[a,b]$, $h=b-a$, $m=\tfrac{a+b}{2}$)}
\renewcommand{\arraystretch}{1.25}
\begin{tabular}{@{}lll@{}}
Regel & Formel & Exaktheit\\\hline
Rechteck & $h\,f(a)$ bzw.\ $h\,f(b)$ & $0$\\
Mittelpunkt & $h\,f(m)$ & $1$\\
Trapez & $\tfrac h2[f(a)+f(b)]$ & $1$\\
Simpson & $\tfrac h6[f(a)+4f(m)+f(b)]$ & $3$\\
$3/8$ & $\tfrac{3h}{8}[f_0+3f_1+3f_2+f_3]$ & $3$\\
Boole & $\tfrac{2h}{45}[7f_0+32f_1+12f_2+32f_3+7f_4]$ & $5$\\
\end{tabular}\\[2pt]
Der \textbf{Mittelpunkt} ist das beste Rechteck: er mittelt die Krümmung (Über-/Unterschätzung
heben sich auf) $\to\mathcal O(h^2)$ statt $\mathcal O(h)$.
\end{satz}

\begin{falle}{Zusammengesetzte Gewichte -- nur die \emph{inneren} Punkte doppeln}
Für $[0,2]$ mit Punkten $f_0,f_1,f_2$ ($h=1$):
\[
  \text{Trapez: } \tfrac h2\bigl[f_0+\mathbf{2}f_1+f_2\bigr],\qquad
  \text{Simpson: } \tfrac h3\bigl[f_0+\mathbf{4}f_1+f_2\bigr].
\]
\textbf{Genau ein} mittlerer Term, mit Gewicht $2$ (Trapez) bzw.\ $4$ (Simpson) -- \emph{nicht}
$f_1+f_2$ zusammenfassen und \emph{nicht} zusätzliche Terme erfinden. Merksatz Simpson:
„$1,4,1$ geteilt durch $3$``.
\end{falle}

\begin{intu}{Warum Simpson Grad 3 exakt schafft}
Simpson legt nur eine \emph{Parabel} (Grad 2) an, ist aber durch \emph{Symmetrie} exakt bis Grad
\textbf{3}: der kubische Fehleranteil hebt sich über das symmetrische Intervall auf.
\end{intu}

\subsection{Interpolation}
\begin{satz}{Lagrange-Polynom}
Durch $n{+}1$ Punkte $(x_i,y_i)$ mit verschiedenen $x_i$ gibt es \emph{genau ein} Polynom vom Grad
$\le n$:
\[
  p(x)=\sum_{i} y_i\,L_i(x),\qquad L_i(x)=\prod_{j\neq i}\frac{x-x_j}{x_i-x_j}.
\]
Jedes $L_i$ ist $1$ in $x_i$ und $0$ in allen anderen Knoten. (Alternativ: Newton-Schema der
dividierten Differenzen -- selbes Polynom, leichter erweiterbar.)
\end{satz}

\begin{falle}{Runge-Phänomen $\Rightarrow$ „stop at Simpson``}
Interpolation \emph{hohen Grades} auf gleichmäßigen Knoten oszilliert stark an den Rändern (Fehler
wächst!). \textbf{Konsequenz}: \emph{nicht} den Polynomgrad erhöhen, sondern \textbf{zusammengesetzte
Regeln niedrigen Grades} verwenden (viele kleine Trapez-/Simpson-Stücke), bzw.\
\textbf{Spline}-Interpolation (stückweise niedriger Grad, glatt, stabil).
\end{falle}

\subsection{Romberg \& Monte-Carlo}
\begin{satz}{Romberg = Richardson auf Trapez}
Trapezwerte für $h,\tfrac h2,\dots$ berechnen und führende Fehlerterme extrapolieren:
$R=\frac{4\,T(h/2)-T(h)}{3}$ (eliminiert den $\mathcal O(h^2)$-Term) $\to$ hohe Genauigkeit aus
billigen Trapezwerten.
\end{satz}

\begin{satz}{Monte-Carlo-Integration}
$\displaystyle \hat I=|\Omega|\,\tfrac1N\sum_{i=1}^N f(X_i)$, Standardfehler
$\mathrm{SE}=\frac{|\Omega|\,\sigma_f}{\sqrt N}\propto N^{-1/2}$. \emph{SE halbieren} $\Rightarrow$
$N\times4$; \emph{zehnteln} $\Rightarrow$ $N\times100$. Die Mittelpunktsregel ist MC mit $N=1$.
\end{satz}

\begin{intu}{Fluch der Dimensionen}
Ein Gitter mit $n$ Punkten pro Achse braucht in $d$ Dimensionen $n^d$ Punkte (exponentiell). Der
MC-Fehler $\propto1/\sqrt N$ ist \emph{dimensionsunabhängig} $\Rightarrow$ in hohen Dimensionen
schlägt Monte-Carlo das Gitter.
\end{intu}

% =====================================================================
\section{Fourier-Analyse \& FFT}

\begin{defi}{Fourier-Transformation}
$F(\omega)=\int_{-\infty}^{\infty} f(t)\,e^{-i\omega t}\,dt$ zerlegt ein Zeitsignal in seine
\textbf{Frequenzanteile} ($|F(\omega)|$ = Stärke der Frequenz $\omega$). Rücktransformation
$f(t)=\frac{1}{2\pi}\int F(\omega)e^{i\omega t}\,d\omega$ setzt das Signal wieder zusammen
(beide invers zueinander).
\end{defi}

\begin{intu}{Euler-Formel \& Interferenz}
$e^{i\omega t}=\cos(\omega t)+i\sin(\omega t)$, also
$\cos(\omega t)=\tfrac12(e^{i\omega t}+e^{-i\omega t})$. Daher liefert ein reeller Kosinus
\emph{zwei} Spektrallinien bei $\pm\omega$. \textbf{Interferenz}: ein Signal ist die \emph{Summe}
mehrerer Schwingungen; $f(t)=A_1\cos(\omega_1 t)+A_2\cos(\omega_2 t)$ zeigt im Spektrum
\textbf{4 Peaks} (bei $\pm\omega_1,\pm\omega_2$, Höhe $\tfrac{A_1}{2},\tfrac{A_2}{2}$).
\end{intu}

\begin{satz}{DFT \& FFT}
\textbf{DFT} (endlich abgetastet): $X_k=\sum_{n=0}^{N-1}x_n\,e^{-i2\pi kn/N}$ -- $N$ diskrete
Frequenz-Bins. Kosten direkt $\mathcal O(N^2)$. \textbf{FFT} spaltet rekursiv in gerade/ungerade
Indizes $\Rightarrow$ $\mathcal O(N\log N)$ (bei $N{=}10^6$ der Unterschied zwischen Sekunden und
Tagen).
\end{satz}

\begin{falle}{Spektrum lesen: Bin-Frequenz \& Nyquist}
Bin $k$ entspricht Frequenz $f_k=k\cdot\frac{f_s}{N}$ ($f_s$ = Abtastrate). \textbf{Nyquist} =
$\frac{f_s}{2}$ ist die höchste eindeutig darstellbare Frequenz. Frequenzen \emph{darüber}
erscheinen fälschlich als niedrigere (\textbf{Aliasing}, Abtasttheorem verletzt). Peaks im
Betragsspektrum $\Leftrightarrow$ periodische Komponenten.
\end{falle}

% =====================================================================
\section{Modelltraining aus ersten Prinzipien}

\begin{intu}{Lineare Regression als Optimierungsproblem}
Modell $\hat y(x)=w x+b$. Trainiert wird durch Minimieren des Verlusts per Gradientenabstieg --
also reine Numerik-Optimierung (Abschnitt 4) auf einer Verlustfunktion.
\end{intu}

\begin{defi}{Mittlerer quadratischer Fehler (MSE) -- mit Quadrat!}
\[
  L(w,b)=\frac1n\sum_{i=1}^n\bigl(\underbrace{w x_i+b}_{\hat y_i}-y_i\bigr)^2 .
\]
Die Abweichungen werden \textbf{quadriert} (sonst heben sich $+$/$-$ auf). Das Quadrat bestraft
große Fehler stärker und macht $L$ glatt/konvex in $(w,b)$.
\end{defi}

\begin{falle}{Erster Loss $=$ Varianz der Zielwerte (Quadrat nicht vergessen!)}
Init $w=0,\ b=\bar y$ sagt überall den Mittelwert voraus, also
\[
  L=\frac1n\sum_i(\bar y-y_i)^2=\mathrm{Var}(y).
\]
\textbf{Varianz = Mittel der \emph{quadrierten} Abweichungen}, \emph{nicht} der Abweichungen selbst
(die ergeben immer $0$). Beispiel $y=(2,4,6)$, $\bar y=4$:
\[
  \mathrm{Var}=\tfrac13\bigl[(-2)^2+0^2+2^2\bigr]=\tfrac13(4+0+4)=\tfrac83\approx2{,}67.
\]
Ein gedruckter Init-Loss $\approx100$ statt $\tfrac83$ deutet auf einen Bug/falsche Initialisierung
(oder Summe statt Mittel); $\approx0$ wäre ebenfalls verdächtig.
\end{falle}

\begin{satz}{Numerik-Brille auf ML-Praktiken}
\begin{itemize}
\item \textbf{Overfit-Sanity-Check}: ein überparametrisiertes Modell \emph{kann} eine kleine
  Teilmenge trivial auf Loss~$0$bringen. Gelingt das \emph{nicht}, liegt der Fehler \emph{nicht} am
  Optimierer, sondern in Modell/Loss/Datenpipeline -- analog: einen Solver an einem Problem mit
  \emph{bekannter exakter Lösung} testen.
\item \textbf{Mini-Batch-SGD} $\leftrightarrow$ \emph{Monte-Carlo}: der Gradient über eine zufällige
  Teilmenge der Größe $B$ ist eine MC-Schätzung; Schätzfehler $\propto 1/\sqrt B$.
\item \textbf{Mehrere Seeds, bester Lauf} $\leftrightarrow$ \emph{Random Restarts} (nicht-konvexe
  Landschaft).
\item \textbf{Quantisierung (niedrige Bit-Präzision)} $\leftrightarrow$ \emph{Präzisionswahl} im
  Gleitkomma: kleinstes Format nehmen, dessen Bereich \& Auflösung genügen (Inferenz verträgt
  weniger Präzision als Training).
\end{itemize}
\end{satz}

\vspace{6pt}
\noindent\textcolor{rulecol}{\rule{\linewidth}{0.4pt}}\\[2pt]
{\small\color{muted} Ende des Lernskripts -- 8 Themenblöcke, von Zahlendarstellung bis Modelltraining.
Die roten \textcolor{red}{Stolperfallen}-Boxen zeigen jeweils direkt die korrekte Form.}

\end{document}
